\section{Related Work}
\label{sec:related}

\textbf{Model Merging and Weight-Space Averaging.} 
Merging pre-trained and fine-tuned models has gained significant traction as a cheap and elegant way to achieve multi-task capabilities~\cite{wortsman2022model,choshen2022fusing}. 
Weight averaging was popularized by Model Soups~\cite{wortsman2022model}, which demonstrated that averaging the weights of multiple fine-tuned models trained with different hyperparameter configurations improves robustness and accuracy. 
To combine models trained on completely different downstream tasks, task arithmetic~\cite{ilharco2022task} introduces the task vector formulation, showing that adding a task vector to base pre-trained weights injects the task's performance. 
However, simple linear addition of weights can suffer from significant parameter interference and cancellation. 
To solve this, Fisher-weighted merging~\cite{matena2022merging} uses Fisher information matrices to weigh individual parameters during averaging. 
TIES-Merging~\cite{yadav2023ties} addresses interference by identifying and resolving sign conflicts, followed by pruning low-magnitude task parameter variations. 
Other methods like Git Re-Basin~\cite{ainsworth2023git} and ZipIt!~\cite{zhang2023zipit} align model symmetries and permutation invariances before merging. 
Our work differs by studying the core \emph{structural resolution} of these merging coefficients rather than introducing new permutation or filtering heuristics.

\textbf{Adaptive Model Merging.} 
Rather than relying on manual coefficient tuning, adaptive model merging frameworks optimize merging scales dynamically. 
AdaMerging~\cite{yang2024adamerging} first proposed optimizing layer-wise coefficients (one coefficient per task per layer) using small unlabeled datasets at test-time, demonstrating that individual layers play distinct roles in multi-task representation routing. 
PolyMerge and SplineMerge~\cite{jung2025poly} expanded on this by parameterizing the layer-wise coefficient profile as a continuous, low-degree polynomial or B-spline of the layer depth, which acts as a low-frequency filter to prevent noisy high-frequency optimization. 
Our GranMerge framework generalizes these approaches by defining a unified structural hierarchy from Level 1 (Global) up to Level 5 (Tensor-wise), allowing us to systematically map out the complete granularity landscape.

\textbf{Test-Time Adaptation and Transductive Overfitting.} 
Our research draws inspiration from test-time adaptation (TTA), where a model is adjusted on unlabeled streams at deployment time. 
TENT~\cite{wang2020tent} minimizes prediction entropy to update normalization statistics, while MEMO~\cite{zhang2021memo} minimizes single-sample prediction entropy across augmented views. 
However, optimizing parameters on small batches is highly prone to transductive overfitting~\cite{mueller2023analyzing,goyal2022test}. 
In multi-task model merging, unconstrained optimization of layer-wise weights on small calibration streams has been shown to cause sacrificial task bias and performance degradation on out-of-distribution samples~\cite{jin2026regcal}. 
RegCalMerge~\cite{jin2026regcal} proposed Class-Capacity Normalization and spatial regularization to mitigate this issue. 
Our study provides the first comprehensive empirical comparison of transductive overfitting across different levels of structural granularity and optimizer families.

\textbf{First-order vs. Zero-order Optimization.} 
Optimizing merging coefficients is traditionally done using first-order gradient descent via backpropagation~\cite{yang2024adamerging,kingma2014adam}. 
However, backpropagation is computationally expensive and requires access to the complete model's internal computational graph. 
Derivative-free zero-order optimization, such as Evolution Strategies (ES)~\cite{rechenberg1973evolutionsstrategie,salimans2017evolution}, offers a black-box alternative. 
ES-based approaches mutate parameters randomly, evaluating only the forward pass loss to guide updates. 
While zero-order methods generally scale poorly with parameter dimensions compared to gradient descent~\cite{hansen2001completely}, they provide unique optimization trajectories and act as robust regularizers. 
We systematically evaluate both optimization paradigms across all granularities.
